% Generated draft: requires mathematical and source review.
\documentclass[11pt]{article}
\usepackage[margin=25mm]{geometry}
\usepackage{fontspec}
\setmainfont{Latin Modern Roman}
\usepackage{amsmath,amssymb,mathtools,unicode-math}
\setmathfont{Latin Modern Math}
\usepackage{enumitem,needspace,xurl,hyperref}
\setlength{\parindent}{0pt}
\setlength{\parskip}{5pt}
\setlength{\emergencystretch}{3em}
\newcommand{\N}{\mathbb{N}}
\newcommand{\Z}{\mathbb{Z}}
\newcommand{\Q}{\mathbb{Q}}
\newcommand{\R}{\mathbb{R}}
\newcommand{\C}{\mathbb{C}}
\newcommand{\F}{\mathbb{F}}
\newcommand{\A}{\mathbb{A}}
\newcommand{\PP}{\mathbb{P}}
\newcommand{\E}{\mathbb{E}}
\newcommand{\eps}{\varepsilon}
\newcommand{\dd}{\,\mathrm{d}}
\begin{document}
\section*{UnsolvedMath ID: 20001311}
\textbf{Source problem code:} catalogue entry
\subsection{Definitions and mathematical statement}
For a convex body \(K\subset\R^n\), \(n\ge2\), its projection body \(\Pi_nK\) is the origin-symmetric convex body with support function
\[
h_{\Pi_nK}(u)=\operatorname{vol}_{n-1}(K|u^\perp),\qquad u\in S^{n-1},
\]
where the bar denotes orthogonal projection. Write \(\Pi^2K=\Pi_n(\Pi_nK)\). W. Weil's recovered Problem 21 asks for the structure of arbitrary convex bodies satisfying
\[
\Pi^2K=cK+x\qquad(c>0,\ x\in\R^n).
\]
The source reports that a polytopal solution is a direct sum of planar polygons in even dimension, or of planar polygons and a segment in odd dimension, and asks what replaces this description without the polytope hypothesis. Here direct sum means a Minkowski sum of factors in complementary linear subspaces, equivalently a linear image of their Cartesian product after choosing coordinates. Central symmetry of the factors is necessary for the centered solution class; the short quoted polytopal description does not state this explicitly.

The supplied aid establishes that every solution is a translate of an origin-symmetric zonoid, a Hausdorff limit of Minkowski sums of centered line segments. If \(z\) is the center of \(K\), then \(x=-cz\), and the centered body \(K-z\) satisfies \(\Pi^2(K-z)=c(K-z)\). Thus arbitrary translation does not add a new shape class.

In dimension two the aid gives the exact identity
\[
\Pi_2^2K=2(K+(-K)).
\]
Consequently the homothetic solutions in the plane are precisely centrally symmetric convex bodies up to translation; the centered multiplier is 4. This includes nonpolygonal solutions and is distinct from classification in higher dimensions.

For a Cartesian product \(K=K_1\times\cdots\times K_r\) in mutually orthogonal complementary subspaces of dimensions \(d_i\), let \(V=\operatorname{vol}_n(K)\), \(V_i=\operatorname{vol}_{d_i}(K_i)\), and \(\sum_i d_i=n\). The supplied product identity is
\[
\Pi_nK=\prod_{i=1}^r\left(\frac{V}{V_i}\Pi_{d_i}K_i\right).
\]
For a one-dimensional factor the convention is \(\Pi_1K_i=[-1,1]\), using zero-dimensional volume one in the support-function formula. If every \(d_i\) is 1 or 2 and all factors are origin symmetric, then
\[
\Pi_n^2K=2^nV^{n-2}K,\qquad
\operatorname{vol}_n(\Pi_nK)=2^nV^{n-1}.
\]
In particular every linear image of such a product is an eigenbody up to homothety. The planar factors may be arbitrary centrally symmetric convex bodies, not only polygons. These examples and identities do not prove that all higher-dimensional solutions are products of this form.

The aid reports that the global classification in dimensions at least three remains open in the literature checked and mentions published local smooth results forcing ellipsoids near a ball. It supplies neither the authors nor the precise topology, neighborhood or regularity hypotheses of those local results, so no unqualified local theorem is asserted here.

The registry labels this entry partially solved; attempt 1 is a partial result, with open status at the run, internal rigor proved and novelty confidence low. The supplied aid is machine generated, not peer reviewed, and requires independent verification before citation.
\subsection{Short English statement}
Weil asks which convex bodies are homothetic to their twice-iterated projection bodies. The aid gives an exact planar answer and explicit products of one- and two-dimensional factors in all dimensions. It reports that these examples do not provide a global classification in dimensions at least three, where the question remains open in the checked literature.
\subsection{Sources}
\begin{itemize}
\item Local UnsolvedMath registry, permanent ID \(20001311\), source code \texttt{catalogue entry}; \texttt{aim-convex-geometry-notes.json}, \texttt{notes[42]}.
\item AIM, \emph{Fourier analytic methods in convex geometry}, recovered Problem 21 (W. Weil): \url{https://aimath.org/WWN/fourierconvex/fourierconvex.pdf}. The verification identifies the lost leading digit in ``1'' and appended page marker ``2''; the permanent ID and source code are unchanged.
\item Supplied research attempt 1 and registry update August 14, 2026. The source's polytopal classification and the aid's claimed published local smooth results have no full bibliographic citations here. The original literature field is empty. No independent PDF or publication inspection is claimed.
\end{itemize}
\end{document}
